% =====================================================================
%  Section 4 -- Experimental Setup
% =====================================================================

\section{Experimental setup}
\label{sec:experimental-setup}

\subsection{Implementation}
\label{sec:setup-impl}

Training is implemented in PyTorch 2.8 with \texttt{segmentation-models-pytorch} 0.4~\cite{iakubovskii2019smp} for UNet decoder construction and \texttt{timm} 1.0~\cite{wightman2019timm} for backbone implementations. Mixed-precision (AMP fp16) training runs on a Kaggle Tesla T4 GPU (16\,GB) with batch size 8 (teacher) or 16 (students). Optimisation uses AdamW with learning rate $3\!\times\!10^{-4}$, weight decay $10^{-4}$, and a cosine schedule. Training runs for 35 epochs (teacher and students alike); empirically, validation IoU plateaus by epoch $\sim 25$. Data augmentation is provided by Albumentations~\cite{buslaev2020albumentations}: random horizontal flip, random rotation $\pm 15^{\circ}$, random brightness/contrast jitter ($\pm 0.2$), and ImageNet normalisation. The random seed is fixed at 42.

ONNX export, INT8 quantization, and CPU/WASM latency benchmarks run locally on an Apple M2~Pro (10-core CPU, 16\,GB unified memory, macOS 26.2); the quantization ablation and the second-host latencies (Section~\ref{sec:setup-latency}) run on an AMD Ryzen~5~5600G desktop (16\,GB, Windows~10). The Kaggle T4 is used for training only.

\subsection{Cross-validation protocol}
\label{sec:setup-cv}

We evaluate folds 0--2 of the five-fold partition \texttt{sklearn.\allowbreak{}model\_selection.\allowbreak{}KFold(n\_splits=5,\allowbreak\,shuffle=True,\allowbreak\,random\_state=42)} (Section~\ref{sec:methods-data} explains the reduction from five folds). The same fold indices are used for every $(\text{model} \times \text{KD configuration})$ combination, so all IoU comparisons in Tables~\ref{tab:t2-segmentation}--\ref{tab:t5-latency} are \emph{paired} across folds.

\subsection{Metrics}
\label{sec:setup-metrics}

We report pixel accuracy, precision, recall, F1-score, and mean intersection-over-union (IoU). For the within-regime zero-shot evaluation (Section~\ref{sec:results-zeroshot}) metrics are computed \emph{per image} and reported as macro statistics, since those test sets mix flood and flood-free scenes: an image with no water in the ground truth and none predicted counts as a correct rejection (IoU 1); over flood-containing images we report the flood-only IoU; and over water-free images we report the false-alarm rate (FAR), the fraction on which the model predicts any water at all. Efficiency metrics are model parameters (M, from \texttt{torchinfo.summary}) and ONNX disk size (MB). Latency is reported as median (p50), 95th percentile (p95), mean, and standard deviation over 200 timed forward passes after 20 warm-up passes; throughput as FPS $= 1000\,/\,\text{p50}_{\text{ms}}$. All latency measurements use single-image batches at $256{\times}256$ input resolution.

\subsection{Statistical validation}
\label{sec:setup-stats}

For each $(\text{model} \times \text{KD configuration})$, we report mean $\pm$ standard deviation across the 3 folds. The paired Wilcoxon signed-rank test is applied to the no-KD vs.~comb-KD contrast on each architecture, using fold-level paired IoU as the unit of observation; with $n=3$ folds its minimum reportable two-sided p-value is 0.25, so we treat it only as a coarse fold-level consistency check. For adequately-powered inference we additionally run a \emph{per-image paired bootstrap}: we pool per-image IoU over the held-out validation images of all three folds ($n=399$ images) and, for each contrast, resample images with replacement ($10{,}000$ replicates) to obtain a 95\% confidence interval on the mean paired IoU difference and a one-sided p-value (\texttt{scripts/stats\_bootstrap.py}; output archived as \texttt{runs/3fold/stats\_bootstrap.json}). Per-image IoU is a biased estimator of the pooled-pixel IoU of Table~\ref{tab:t2-segmentation} (small-flood images weigh as much as large ones), but the bias is image-specific and cancels in the \emph{paired} difference, which is the estimand. This bootstrap, not the three-fold Wilcoxon, supports the significance statements in Section~\ref{sec:results-kd}.

\subsection{Per-platform latency measurement}
\label{sec:setup-latency}

\begin{itemize}
\item \textbf{Apple M2~Pro / ONNX Runtime 1.22.1 CPU EP:} \texttt{scripts/bench\_m2.py}. Default ONNX Runtime CPU threading (the intra-op and inter-op thread counts are left at their defaults; the M2~Pro exposes up to 10 intra-op threads), matching realistic laptop deployment where all cores are available. 20 warm-up $+$ 200 timed runs.
\item \textbf{Browser WebAssembly SIMD / ONNX Runtime Web 1.17 from Node.js~25:} \texttt{bench.mjs} (supplementary). WASM execution provider with SIMD enabled and \texttt{numThreads} set to 1. 20 warm-up $+$ 200 timed runs. Equivalent in-browser numbers (Chrome 130, macOS 26.2) were verified to fall within $\pm 10\%$ in pilot runs.
\item \textbf{Second host, AMD Ryzen~5~5600G (6 cores, Windows~10):} \texttt{scripts/bench\_ort\_cpu.py} (ONNX Runtime 1.22.1 CPU EP, default threading) and \texttt{scripts/bench\_wasm\_node.mjs} (ONNX Runtime Web 1.22.0, WASM SIMD, one thread, Node.js~22), same warm-up and run counts. Timed on the fold-0 FP32 and INT8 artifacts produced by the quantization ablation (Section~\ref{sec:results-quant}), which are opset-17 exports; FP32 and INT8 are therefore compared at the same opset on this host.
\end{itemize}

All benchmarks consume the identical \texttt{.onnx} artifact; no platform-specific graph rewrites are performed.
